\chapter{الفضاءات المتجهية المعيارية}\label{ch:b2:nvs}

حين يكون \omterm{def:b2:metric:def}{الفضاء المتري} فضاءً متجهيًا وتأتي المسافة من
\omterm{def:b2:nvs:norm}{معيار}، تبدأ الطوبولوجيا \omterm{def:b2:structures:algebra}{والجبر} الخطي بالتفاعل: فالتطبيقات الخطية
تكون متصلة بالضبط حين تكون محدودة على كرة الوحدة،
ويفرض البعد المنتهي تكافؤ كل المعايير، ويحوّل التمام
المتسلسلات المتقاربة بإطلاق إلى متسلسلات متقاربة. والفاصل
بين البعد المنتهي وغير المنتهي --- المتبلور في مبرهنة ريس ---
هو أعمق درس في الفصل.

وفي كل ما يلي، $E, F$ فضاءات متجهية على $K = \R$ أو $\C$.

\section{المعايير}

\begin{definition}\label{def:b2:nvs:norm}
\emph{المعيار}\index{معيار} على $E$ تطبيق $\norm{\,\cdot\,} \colon E
\to \R_+$ يحقق، من أجل كل $x, y \in E$ و$\lambda \in K$:
\[
\norm x = 0 \iff x = 0,
\qquad
\norm{\lambda x} = \abs\lambda\,\norm x,
\qquad
\norm{x + y} \leq \norm x + \norm y .
\]
عندئذٍ يكون $d(x, y) = \norm{x - y}$ مسافة، وينطبق كل ما في
\cref{ch:b2:metric}. وتجعل متراجحة المثلث المعكوسة
$\bigl|\norm x - \norm y\bigr| \leq \norm{x - y}$ المعيارَ نفسه
\omterm{def:b2:metric:continuity}{ليبشيتزيًا} بالثابت $1$؛ والجمع والضرب السلّمي
متصلان (بالتقديرات $\norm{(x + y) - (x' + y')} \leq \norm{x -
x'} + \norm{y - y'}$، وغيرها).
\end{definition}

\begin{example}\label{ex:b2:nvs:examples}
على $K^n$:
\[
\norm{x}_1 = \sum_i \abs{x_i},
\qquad
\norm{x}_2 = \Bigl(\sum_i \abs{x_i}^2\Bigr)^{1/2},
\qquad
\norm{x}_\infty = \max_i \abs{x_i}
\]
(و$\norm\cdot_2$ \omterm{def:b2:nvs:norm}{معيار} بمتراجحة كوشي--شوارتز، مجلد السنة الأولى). وعلى
$C(\intcc{a}{b})$:
\[
\norm f_\infty = \sup \abs f,
\qquad
\norm f_1 = \int_a^b \abs f,
\qquad
\norm f_2 = \Bigl(\int_a^b \abs f^2\Bigr)^{1/2},
\]
والأخيران معياران بفضل الإيجابية التامة
للتكامل ومتراجحة كوشي--شوارتز التكاملية (مجلد السنة الأولى). وعلى المصفوفات:
أي \omterm{def:b2:nvs:norm}{معيار} على $\mathcal M_n(K) \simeq K^{n^2}$؛ ومعايير المؤثرات
أدناه هي المهمة بنيويًا.
\end{example}

\begin{definition}[المعايير المتكافئة]\label{def:b2:nvs:equivalent}
يكون معياران $N_1, N_2$ على $E$ \emph{متكافئين}\index{معايير
متكافئة} إذا وُجد ثابتان $c, C > 0$ يحققان
\[
c\,N_1 \leq N_2 \leq C\, N_1 .
\]
وللمعايير المتكافئة المفتوحات نفسها، والمتتاليات المتقاربة
والكوشية نفسها، والأجزاء المتراصة والتامة نفسها: أي التحليل نفسه.
\end{definition}

\begin{example}[عدم التكافؤ في البعد غير المنتهي]\label{ex:b2:nvs:nonequivalent}
على $C(\intcc{0}{1})$: $\norm f_1 \leq \norm f_\infty$ دائمًا، لكن
لا يصحّ أي حدّ معاكس: فالدالة $f_n(x) = x^n$ لها $\norm{f_n}_\infty = 1$
و$\norm{f_n}_1 = \frac{1}{n+1} \to 0$. ومنه $f_n \to 0$ من أجل
$\norm\cdot_1$ لا من أجل $\norm\cdot_\infty$: فالمعياران
يختلفان في التقارب نفسه.
\end{example}

\begin{example}[ثوابت صريحة في البعد $n$]\label{ex:b2:nvs:constants}
على $K^n$ تتكافأ المعايير الكلاسيكية الثلاثة بثوابت
مثلى:
\[
\norm x_\infty \leq \norm x_2 \leq \norm x_1
\leq \sqrt n\,\norm x_2 \leq n\,\norm x_\infty ,
\]
والحدّ الأوسط $\norm x_1 \leq \sqrt n\norm x_2$ آتٍ من
متراجحة كوشي--شوارتز في مقابل المتجهة ذات الواحدات. والمتجهات الحدّية:
$e_1$ تجعل المتراجحتين الأوليين مساواتين، و$(1, 1,
\dots, 1)$ تجعل الأخيرتين كذلك. ويظهر البعد $n$ بوضوح في
الثوابت --- وهو البذرة الكمّية للفشل في البعد غير
المنتهي: فحين $n \to \infty$ لا يبقى أي ثابت منتظم،
وهذا بالضبط ما يُبرزه \cref{ex:b2:nvs:nonequivalent} على
فضاءات الدوال.
\end{example}

\begin{omfigure}
\begin{tikzpicture}[scale=1.5]
  \draw[->] (-1.45,0) -- (1.5,0) node[below, font=\small] {$x$};
  \draw[->] (0,-1.45) -- (0,1.5) node[left, font=\small] {$y$};
  \draw[omProp, very thick]
    (1,1) -- (-1,1) -- (-1,-1) -- (1,-1) -- cycle;
  \draw[omThm, very thick] (0,0) circle (1);
  \draw[omDef, very thick]
    (1,0) -- (0,1) -- (-1,0) -- (0,-1) -- cycle;
  \node[omDef, font=\small] at (0.36,0.28)
    {$\norm\cdot_1$};
  \node[omThm, font=\small] at (0.55,0.92)
    {$\norm\cdot_2$};
  \node[omProp, font=\small] at (1.24,1.13)
    {$\norm\cdot_\infty$};
\end{tikzpicture}

{\small كرات الوحدة للمعايير الكلاسيكية الثلاثة في $\R^2$،
متداخلة كما تقتضي متراجحات \cref{ex:b2:nvs:constants}:
فكلما صغرت الكرة كبر \omterm{def:b2:nvs:norm}{المعيار}. والاستدارة مهمة: فالأضلاع
المستقيمة للمعيّن والمربّع هي بالضبط مواضع فشل
التحدب التام المستغَلة في مسألة نهاية الأسبوع في \cref{ch:b2:realfun}
وفي مسألة هذا الفصل (السؤال 4).}
\end{omfigure}

\section{التطبيقات الخطية المتصلة}

\begin{theorem}[توصيف]\label{thm:b2:nvs:continuouslinear}
من أجل تطبيق خطي $u \colon E \to F$ بين فضاءين معياريين، تتكافأ
الأقوال التالية:
\begin{enumerate}
  \item $u$ \omterm{def:b2:metric:continuity}{متصل}؛
  \item $u$ \omterm{def:b2:metric:continuity}{متصل} عند $0$؛
  \item $u$ محدود على كرة الوحدة المغلقة: $\sup_{\norm x \leq
        1} \norm{u(x)} < \infty$؛
  \item يوجد $C \geq 0$ يحقق $\norm{u(x)} \leq C \norm x$ من أجل
        كل $x$؛
  \item $u$ \omterm{def:b2:metric:continuity}{ليبشيتزي}.
\end{enumerate}
وأصغر عدد $C$ كهذا هو \emph{معيار المؤثر}\index{معيار المؤثر}
$\vertiii{u} = \sup_{\norm x \leq 1}\norm{u(x)} = \sup_{x \neq 0}
\frac{\norm{u(x)}}{\norm x}$؛ وهو يجعل فضاء التطبيقات الخطية المتصلة $\mathcal{L}_c(E,
F)$ فضاءً معياريًا، مع
\[
\vertiii{v \circ u} \leq \vertiii v\, \vertiii u .
\]
\end{theorem}

\begin{proof}
(1 $\Rightarrow$ 2) بديهي. (2 $\Rightarrow$ 3): يعطي الاتصال عند $0$
مع $\varepsilon = 1$ عددًا $\delta$ يحقق $\norm x \leq \delta
\Rightarrow \norm{u(x)} \leq 1$؛ ثم تُنزل التجانسية أي $x$
يحقق $\norm x \leq 1$ إلى تلك الكرة وتعيده:
\[
\norm{u(x)} = \frac1\delta\,\norm{u(\delta x)} \leq
\frac1\delta ,
\]
لأن $\norm{\delta x} \leq \delta$. (3 $\Rightarrow$ 4):
من أجل $x \neq 0$، نطبّق الحدّ على $\frac{x}{\norm x}$. (4
$\Rightarrow$ 5): $\norm{u(x) - u(y)} = \norm{u(x - y)} \leq
C\norm{x - y}$. (5 $\Rightarrow$ 1) معلوم.

بديهيات \omterm{def:b2:nvs:norm}{المعيار} من أجل $\vertiii\cdot$: التجانسية والفصل
واضحان (فالشرط $\vertiii u = 0$ يفرض $u = 0$ على الكرة، ومن ثم في
كل مكان)؛ ومتراجحة المثلث من $\norm{(u + v)(x)} \leq
\norm{u(x)} + \norm{v(x)}$. والضربية التحتية: $\norm{v(u(x))}
\leq \vertiii v\,\norm{u(x)} \leq \vertiii v \vertiii u \norm x$.
\end{proof}

\begin{example}\label{ex:b2:nvs:operatornorms}
على $\bigl(C(\intcc{0}{1}), \norm\cdot_\infty\bigr)$: للتقييم
$f \mapsto f(0)$ \omterm{def:b2:nvs:norm}{معيار} مؤثر $1$؛ وللتكامل $f \mapsto
\int_0^1 f$ \omterm{def:b2:nvs:norm}{معيار} $1$؛ وللتطبيق $f \mapsto \int_0^1 t f(t)\dd t$
\omterm{def:b2:nvs:norm}{معيار} $\int_0^1 t\,\dd t = \frac12$ (والحدّ الأعلى بمتراجحة المثلث
للتكاملات؛ وهو مبلوغ عند $f \equiv 1$). أما الاشتقاق،
من $(C^1, \norm\cdot_\infty)$ إلى $(C^0,
\norm\cdot_\infty)$، \emph{فليس} \omterm{def:b2:metric:continuity}{متصلًا}: إذ إن $\norm{\sin(nx)}_\infty
= 1$ في حين أن \omterm{def:b2:nvs:norm}{معيار} النهاية العليا للمشتق $n$. فالخطية لا تستلزم
الاتصال في البعد غير المنتهي.
\end{example}

\begin{example}[معياران، وحكمان على متتالية واحدة]\label{ex:b2:nvs:sqrtnxn}
على $C(\intcc01)$، لتكن $g_n(x) = \sqrt{n}\,x^n$. عندئذٍ
\[
\norm{g_n}_1 = \frac{\sqrt n}{n + 1} \longrightarrow 0,
\qquad
\norm{g_n}_2^2 = \frac{n}{2n + 1} \longrightarrow \frac12,
\qquad
\norm{g_n}_\infty = \sqrt n \longrightarrow \infty :
\]
متتالية واحدة، وثلاثة معايير، وثلاثة سلوكات --- تقارب إلى
الصفر، وعدم تقارب (فالمعايير تستقر عند $\frac1{\sqrt2}$ لكن
النهاية النقطية $0$)، وانفجار. وتركّز الكتلة
قرب $x = 1$ غير مرئي \omterm{def:b2:nvs:norm}{للمعيار} $\norm\cdot_1$، ومرئي نصفًا \omterm{def:b2:nvs:norm}{للمعيار}
$\norm\cdot_2$، ومهيمن من أجل $\norm\cdot_\infty$. وفي البعد
غير المنتهي، ليس السؤال ``هل تتقارب؟'' سؤالًا عن
متتالية: بل هو سؤال عن متتالية \emph{\omterm{def:b2:nvs:norm}{ومعيار}}.
\end{example}

\begin{method}[كيف نحسب معيار مؤثر]\label{met:b2:nvs:opnorm}
دائمًا بحركتين. \emph{الحدّ الأعلى:} قدّر
$\norm{u(x)}$ بالمقدار $C\norm x$ باستعمال متراجحات المثلث
أو كوشي--شوارتز أو حدود التكاملات --- وهذا يبرهن على $\vertiii u
\leq C$. \emph{الشاهد:} أبرِز إما عنصرًا بعينه $x_0 \neq 0$
يحقق $\norm{u(x_0)} = C\norm{x_0}$ (فيكون الحدّ مبلوغًا)، وإما
متتالية من متجهات الوحدة $x_n$ تحقق $\norm{u(x_n)} \to C$ (فيكون
الحدّ مقارَبًا). والحركتان إلزاميتان: فالحدّ الأعلى
وحده لا يعطي إلا $\vertiii u \leq C$، والشاهد وحده لا يعطي إلا
$\vertiii u \geq C$. وفي البعد غير المنتهي قد يضطر الشاهد
أن يكون متتالية --- إذ لا يلزم أن تكون النهاية العليا مبلوغة
(\cref{exo:b2:nvs:8}).
\end{method}

\begin{example}[المؤثرات القطرية ترى كل المعايير سواءً]\label{ex:b2:nvs:diagonalnorm}
من أجل $D = \operatorname{diag}(d_1, \dots, d_n)$ على $K^n$ \emph{بأيٍّ} من
المعايير $\norm\cdot_1, \norm\cdot_2,
\norm\cdot_\infty$: من $\abs{d_ix_i} \leq
\bigl(\max_j\abs{d_j}\bigr)\abs{x_i}$ إحداثيةً إحداثية، نجد
$\norm{Dx} \leq \max_j\abs{d_j}\,\norm x$؛ ويبلغه $x = e_{j_0}$
(بدليل معظِّم). ومنه $\vertiii D =
\max_j\abs{d_j}$ في الحالات الثلاث: فمن أجل التطبيقات القطرية، تحكي كل
المعايير المعقولة الحكاية نفسها، وهي أكبر معامل تمديد.
وكل ما هو صعب في معايير المؤثرات يتعلق بالسلوك
\emph{غير} القطري --- ولهذا تعمل المعايير المكيَّفة في
مسألة نهاية الأسبوع في هذا الفصل (السؤال 22) بأن
تجبر المصفوفة على أن تصير قطرية أولًا.
\end{example}

\begin{example}[مجاميع الأعمدة: توأم $1$ للمعيار \cref{exo:b2:nvs:4}]\label{ex:b2:nvs:columnsums}
على $(\R^n, \norm\cdot_1)$، \omterm{def:b2:nvs:norm}{معيار} مؤثر مصفوفة $A$ هو
أكبر مجموع مطلق \emph{لعمود}. ونشغّل الطريقة: من أجل
$\norm x_1 \leq 1$،
\[
\norm{Ax}_1 = \sum_i\Bigl|\sum_j a_{ij}x_j\Bigr|
\leq \sum_j \abs{x_j}\sum_i\abs{a_{ij}}
\leq \Bigl(\max_j\sum_i\abs{a_{ij}}\Bigr)\norm x_1 ,
\]
ويُبلغ الحدّ عند $x = e_{j_0}$ من أجل عمود معظِّم
$j_0$ --- وهو أنظف شاهد يمكن تصوّره. ومن أجل $A =
\begin{psmallmatrix}1 & -2\\ 3 & 1\end{psmallmatrix}$:
$\vertiii A_1 = \max(1 + 3,\ 2 + 1) = 4$، في حين أن
$\vertiii A_\infty = 4$ أيضًا (بالأسطر) --- وهي مصادفة هنا لا
قانون: فبدّل عناصر المصفوفة بغير تناظر ويفترق
المعياران. الأسطر من أجل $\norm\cdot_\infty$، والأعمدة من أجل
$\norm\cdot_1$: والقاعدة الذهنية أن متجهات الوحدة لكل \omterm{def:b2:nvs:norm}{معيار}
(أنماط الإشارات، وأما متجهات الأساس فللآخر) تلتقط المجاميع
الموافقة.
\end{example}

\begin{proposition}[التطبيقات ثنائية الخطية]\label{prop:b2:nvs:bilinear}
يكون التطبيق ثنائي الخطية $b \colon E \times F \to G$ \omterm{def:b2:metric:continuity}{متصلًا} إذا وفقط إذا كان
$\norm{b(x,y)} \leq C\norm x\,\norm y$ من أجل $C$ ما؛ ويكون عندئذٍ
\omterm{def:b2:metric:continuity}{ليبشيتزيًا} على المجموعات المحدودة. (بنمط البرهان نفسه؛ والجداء
$(u, v) \mapsto v \circ u$ وضرب المصفوفات هما المثالان
الأساسيان.)
\end{proposition}

\begin{proof}
إذا صحّ الحدّ:
\[
b(x,y) - b(x_0,y_0) = b(x - x_0,\, y) + b(x_0,\, y - y_0),
\]
ومنه $\norm{b(x,y) - b(x_0,y_0)} \leq C\norm{x - x_0}\norm y +
C\norm{x_0}\norm{y - y_0}$: أي الاتصال عند $(x_0, y_0)$، وحدّ
\omterm{def:b2:metric:continuity}{ليبشيتزي} حيث $\norm x, \norm y \leq R$. وعكسيًا، يعطي
الاتصال عند $(0,0)$ عددًا $\delta$ يحقق $\norm{b(x,y)} \leq 1$ على
$\norm x, \norm y \leq \delta$؛ ثم نغيّر سلّم المتغيرين.
\end{proof}

\section{البعد المنتهي}

\begin{theorem}[تكافؤ المعايير في البعد المنتهي]\label{thm:b2:nvs:finitedim}
على فضاء منتهي البعد، \emph{كل \omterm{def:b2:nvs:equivalent}{المعايير
متكافئة}}.\index{تكافؤ المعايير} ومن ثم، في البعد
المنتهي: يكون التقارب والانفتاح والتراص والتمام مفاهيم
لا تتعلق \omterm{def:b2:nvs:norm}{بالمعيار}؛ ويكون المتراص $=$ مغلقًا ومحدودًا؛
ويكون الفضاء \omterm{def:b2:metric:complete}{تامًا}؛ ويكون كل تطبيق خطي (أو متعدد الخطية)
\emph{منطلق} من فضاء منتهي البعد \omterm{def:b2:metric:continuity}{متصلًا}.
\end{theorem}

\begin{proof}
نثبّت أساسًا ونماهي $E \simeq K^n$؛ ويكفي أن نقارن أي
\omterm{def:b2:nvs:norm}{معيار} $N$ \omterm{def:b2:nvs:norm}{بالمعيار} $\norm\cdot_\infty$.

\emph{أحد الاتجاهين جبر:} $N(x) = N(\sum x_i e_i) \leq
\sum \abs{x_i} N(e_i) \leq C \norm x_\infty$ مع $C = \sum N(e_i)$.
وهذا يبيّن أيضًا أن $N$ \omterm{def:b2:metric:continuity}{متصل} على $(K^n, \norm\cdot_\infty)$
(فهو \omterm{def:b2:metric:continuity}{ليبشيتزي} بالثابت $C$: $\abs{N(x) - N(y)} \leq N(x - y)$).

\emph{والآخر طوبولوجيا:} الكرة الواحدية $S = \{x :
\norm{x}_\infty = 1\}$ مغلقة ومحدودة في $(K^n,
\norm\cdot_\infty)$، ومن ثم متراصة
(\cref{thm:b2:metric:compactprops}~(2)، وهي صالحة من أجل $\C^n \simeq
\R^{2n}$). وتبلغ الدالة المتصلة $N$ حدّها الأدنى $c$ على
$S$؛ و$c > 0$ لأن $N$ لا ينعدم إلا عند $0 \notin S$. وتنشر التجانسية
الحدَّ: $N(x) \geq c \norm{x}_\infty$ من أجل كل $x$.

النتائج: تؤول كل الأقوال إلى الفضاء $(K^n,
\norm\cdot_\infty)$، حيث هي معلومة
(\cref{thm:b2:metric:rncomplete}،
\cref{thm:b2:metric:compactprops})؛ ويحقق التطبيق الخطي $u$ المنطلق من
فضاء منتهي البعد $E$ العلاقةَ $\norm{u(x)} \leq \sum\abs{x_i}
\norm{u(e_i)} \leq C'\norm{x}_\infty$: أي الحدّ (4) في
\cref{thm:b2:nvs:continuouslinear}.
\end{proof}

\begin{corollary}\label{cor:b2:nvs:closedsubspace}
كل فضاء جزئي منتهي البعد من أي فضاء معياري مغلق.
\end{corollary}

\begin{proof}
هو تام من أجل \omterm{def:b2:nvs:norm}{المعيار} المستحثّ
(\cref{thm:b2:nvs:finitedim})، والأجزاء التامة مغلقة
(\cref{def:b2:metric:complete}).
\end{proof}

\begin{example}[أفضل تقريب محسوب بالتناظر]\label{ex:b2:nvs:bestaffine}
في $\bigl(C(\intcc{-1}{1}), \norm\cdot_\infty\bigr)$، كم تبعد
$f(x) = \abs x$ عن الفضاء الجزئي (المغلق ذي البعد اثنين) المكوَّن
من الدوال الأفينية $a + bx$؟ بالتناظر، لا يغيّر تعويض $a + bx$ بالمقدار
$a - bx$ المقدارَ $\norm{f - (a \pm bx)}_\infty$، وتفي النقطة
الوسطى $a$ بالغرض على الأقل (بمتراجحة المثلث على
المتوسط): فيكفي النظر في الثوابت. ومن أجل ثابت
$a$:
\[
\norm{\abs x - a}_\infty = \max\,(1 - a,\ a)
\geq \frac12 ,
\]
وهو أصغري عند $a = \frac12$: فالمسافة $\frac12$، وهي مبلوغة
بالثابت $\frac12$. ولاحظ منحني الخطأ $\abs x -
\frac12$: فهو يبلغ $\pm\frac12$ بالتناوب عند $x = -1, 0,
1$ --- أي ثلاثة حدود قصوى متناوبة الإشارة من أجل أفضل
تقريب من عائلة ذات وسيطين. وليس نمط
\emph{التذبذب المتساوي} هذا مصادفةً؛ بل هو
بصمة الأمثلية التي تحوّلها مسألة نهاية الأسبوع في هذا الفصل
إلى مبرهنة تشيبيشيف.
\end{example}

\begin{example}[الفضاءات الجزئية المغلقة في مقابل الكثيفة]\label{ex:b2:nvs:closeddense}
في $E = \bigl(C(\intcc{0}{1}), \norm\cdot_\infty\bigr)$: كل فضاء
$\R_n[X]$ (كثيرات الحدود من الدرجة $\leq n$، مقصورة على
$\intcc01$) فضاء جزئي منتهي البعد، ومن ثم \emph{مغلق}
--- فالنهاية المنتظمة لكثيرات حدود من الدرجة $\leq n$ كذلك. لكن
اتحادها كلها $\R[X]$ \emph{كثيف} في $E$ (وهي
مبرهنة فايرشتراس في التقريب، المبرهَن عليها في
\cref{ch:b2:funcseq})، والفضاءات الجزئية الفعلية الكثيفة
غير مغلقة إلى أقصى حدّ. والعبرة: انغلاق الفضاءات الجزئية
امتياز للبعد المنتهي؛ وتكديس طوابق مغلقة قد يبني
ناطحة سحاب كثيفة.
\end{example}

\begin{theorem}[ريس]\label{thm:b2:nvs:riesz}
تكون كرة الوحدة المغلقة لفضاء معياري $E$ متراصة \emph{إذا وفقط
إذا} كان $\dim E < \infty$.\index{مبرهنة ريس}
\end{theorem}

\begin{proof}
البعد المنتهي: يكفي أنها مغلقة ومحدودة
(\cref{thm:b2:nvs:finitedim}).

وعكسيًا، لنفترض $\dim E = \infty$. \emph{مبرهنة ريس المساعدة:} من أجل
كل فضاء جزئي مغلق فعلي $F \subsetneq E$ وكل $\varepsilon \in
\intoo{0}{1}$، توجد متجهة واحدية $x$ تحقق $d(x, F) \geq 1 -
\varepsilon$. البرهان: نختار $y \notin F$، ونضع $\delta = d(y, F) > 0$
(وهي موجبة لأن $F$ مغلق)، ونختار $f \in F$ تحقق $\norm{y - f} \leq
\frac{\delta}{1 - \varepsilon}$، ونضع $x = \frac{y - f}{\norm{y -
f}}$: فمن أجل أي $g \in F$،
\[
\norm{x - g} = \frac{\norm{y - (f + \norm{y-f}\,g)}}{\norm{y - f}}
\geq \frac{\delta}{\norm{y-f}} \geq 1 - \varepsilon ,
\]
لأن البسط مسافة من $y$ إلى نقطة من $F$.

ثم نبني متجهات واحدية $x_1, x_2, \dots$ بالتراجع: فالفضاء $F_k =
\operatorname{Vect}(x_1, \dots, x_k)$ منتهي البعد، ومن ثم
مغلق (\cref{cor:b2:nvs:closedsubspace}) وفعلي؛ وتعطي مبرهنة ريس المساعدة
مع $\varepsilon = \frac12$ متجهةً واحدية $x_{k+1}$ تحقق $d(x_{k+1},
F_k) \geq \frac12$. وتحقق المتتالية $\norm{x_p - x_q} \geq
\frac12$ من أجل $p \neq q$: فلا متتالية جزئية متقاربة --- أي إن كرة الوحدة
ليست متراصة.
\end{proof}

\begin{example}[ريس بوصفه كاشفًا للبعد]\label{ex:b2:nvs:rieszdetector}
هل $C(\intcc01)$ منتهي البعد؟ يجيب ريس دون
إبراز أي عائلة حرة غير منتهية صريحة: فالمتتالية
$f_n(x) = x^n$ تقع في كرة الوحدة المغلقة وتحقق، من أجل
$m > n$، $\norm{f_n - f_m}_\infty \geq f_n(x_0) - f_m(x_0) > 0$
عند نقاط مناسبة --- وهو ما يُقدَّر بدقة في مسألة نهاية
الأسبوع في هذا الفصل (السؤال 16)، حيث تبقى متتالية جزئية على
مسافة متبادلة $\geq \frac14$. فلا متتالية جزئية متقاربة، ومن ثم
فالكرة غير متراصة، ومنه $\dim C(\intcc01) = \infty$ حسب
\cref{thm:b2:nvs:riesz}. وتراص كرة الوحدة انقسام
تامّ: فهو يصحّ في البعد المنتهي ويفشل في
البعد غير المنتهي، بلا أرض وسطى --- فالهندسة وحدها
تقرأ نوع البعد.
\end{example}

\section{فضاءات باناخ}

\begin{definition}\label{def:b2:nvs:banach}
\emph{فضاء باناخ}\index{فضاء باناخ} فضاء معياري تام.
أمثلة: كل فضاء معياري منتهي البعد
(\cref{thm:b2:nvs:finitedim})؛ و$\bigl(C(\intcc{a}{b}),
\norm\cdot_\infty\bigr)$ (\cref{thm:b2:metric:rncomplete})؛
و$\mathcal{L}_c(E, F)$ من أجل $F$ فضاءَ باناخ (بنمط البرهان نفسه المستعمل من أجل
الدوال المتصلة). ومثال مضاد: $\bigl(C(\intcc{0}{1}),
\norm\cdot_1\bigr)$ (\cref{exo:b2:nvs:7}).
\end{definition}

\begin{example}[معيار مؤثر التكامل]\label{ex:b2:nvs:integrationnorm}
على $\bigl(C(\intcc01), \norm\cdot_\infty\bigr)$، ليكن $T(f)(x) =
\int_0^x f(t)\,\dd t$ (وهو تشاكل ذاتي: فالدالة $T(f)$ متصلة).
ونشغّل \cref{met:b2:nvs:opnorm}. الحدّ الأعلى:
\[
\abs{T(f)(x)} \leq \int_0^x\abs f \leq x\,\norm f_\infty \leq
\norm f_\infty ,
\]
ومنه $\vertiii T \leq 1$. والشاهد: يعطي $f \equiv 1$ أن $T(f)(x) =
x$ و$\norm{T(f)}_\infty = 1 = \norm f_\infty$: فهو مبلوغ،
و$\vertiii T = 1$. لكن لاحظ أن $\vertiii{T^2} = \frac12 \neq
\vertiii T^2$: فعلًا، $T^2(f)(x) = \int_0^x(x - t)f(t)\dd t$ له
$\abs{T^2(f)(x)} \leq \frac{x^2}2\norm f_\infty$، وهو مبلوغ من جديد
عند $f \equiv 1$؛ وعمومًا $\vertiii{T^n} = \frac1{n!}$ ---
فالحدّ الضربي التحتي $\vertiii T^n = 1$ يخطئ بمقدار
عاملي. وهذه بالضبط هي الظاهرة التي تحوّلها حيلة التكرار في
مسألة نهاية الأسبوع في \cref{ch:b2:metric} إلى قابلية حلّ
شمولية للمعادلات التفاضلية الخطية.
\end{example}

\begin{theorem}[التقارب المطلق في فضاءات باناخ]\label{thm:b2:nvs:absoluteconvergence}
في \omterm{def:b2:nvs:banach}{فضاء باناخ}، إذا كان $\sum \norm{u_n} < \infty$ فإن $\sum u_n$
تتقارب، ويكون $\norm{\sum u_n} \leq \sum\norm{u_n}$. (وتُعالَج
نظرية المتسلسلات في \omterm{def:b2:nvs:norm}{الفضاءات المعيارية} كاملةً في \cref{ch:b2:series}.)
\end{theorem}

\begin{proof}
المجاميع الجزئية $S_N$: من أجل $q > p$، $\norm{S_q - S_p} \leq
\sum_{n=p+1}^{q}\norm{u_n}$، وهذا يؤول إلى $0$ (بمحك كوشي
للمتسلسلة الحقيقية للمعايير): ومن ثم فإن $(S_N)$ كوشي، فمتقاربة.
وتنتقل المتراجحة إلى النهاية من متراجحة المثلث
المنتهية.
\end{proof}

\begin{example}[الأسّي المصفوفي، أول لقاء]\label{ex:b2:nvs:matrixexp}
\emph{الأسّي المصفوفي}\index{أسّي مصفوفي}:
$\mathcal{M}_n(K)$ مع أي \omterm{def:b2:nvs:norm}{معيار} ضربي تحتي ($\vertiii{AB}
\leq \vertiii A \vertiii B$) \omterm{def:b2:nvs:banach}{فضاء باناخ} (لأنه منتهي البعد). عندئذٍ، من أجل
كل $A$،
\[
\eu^A = \sum_{k=0}^{\infty} \frac{A^k}{k!}
\]
تتقارب بإطلاق (لأن $\vertiii{A^k/k!} \leq \vertiii A^k /k!$،
وهي قابلة للجمع): فهي معرَّفة جيدًا. ويستغلّ \cref{ch:b2:diffeq} ذلك
منهجيًا.
\end{example}

\begin{example}[متسلسلة نويمان تنتهي]\label{ex:b2:nvs:neumannfinite}
من أجل $A = \begin{psmallmatrix}0 & \frac12\\ 0 &
0\end{psmallmatrix}$: $\vertiii A < 1$ في أي \omterm{def:b2:nvs:norm}{معيار} مؤثر
مبني على معايير \cref{ex:b2:nvs:examples}، و$A^2 = 0$، ومن ثم
تنهار المتسلسلة الهندسية:
\[
(I - A)^{-1} = \sum_{k \geq 0} A^k = I + A =
\begin{pmatrix}1 & \tfrac12\\ 0 & 1\end{pmatrix},
\]
وهو ما يتحقق منه $(I - A)(I + A) = I - A^2 = I$. فانعدام القوى
يقتطع المتسلسلة تمامًا كما اقتطع الأسّي في
\cref{ch:b2:reduction}؛ ويعاير المثال
التوقعات: فمقلوب نويمان متسلسلة غير منتهية
عمومًا، وكثيرُ حدود بالضبط حين يكون الاضطراب
عديم القوى، والخطأ بعد $N$ حدًّا محدود دائمًا
بالذيل الهندسي $\vertiii A^{N+1}/(1 - \vertiii A)$.
\end{example}

\begin{example}[أسّي مولّد دوران]\label{ex:b2:nvs:rotationexp}
لتكن $A = \begin{psmallmatrix}0 & -\theta\\ \theta &
0\end{psmallmatrix}$. عندئذٍ $A^2 = -\theta^2 I$، ومن ثم تدور
القوى بدور قدره أربعة، وتنقسم المتسلسلة إلى جزأين زوجي
وفردي:
\[
\eu^{A} = \sum_{k}\frac{A^k}{k!}
= \Bigl(\sum_{j}\frac{(-1)^j\theta^{2j}}{(2j)!}\Bigr) I
+ \Bigl(\sum_{j}\frac{(-1)^j\theta^{2j+1}}{(2j+1)!}\Bigr)
\frac{A}{\theta}
= \begin{pmatrix}
\cos\theta & -\sin\theta\\
\sin\theta & \cos\theta
\end{pmatrix},
\]
وكل إعادة ترتيب مرخَّصة بالتقارب المطلق. فأسّي
مولّد ضدّ متماثل دوران ---
محسوبٌ هنا من المتسلسلة وحدها، قبل ثلاثة فصول من
شرح المعادلة التفاضلية $x' = Ax$ (\cref{ch:b2:diffeq})
\emph{لماذا}: فالمقدار $\eu^{tA}$ حركة دائرية منتظمة. والفكرة
الختامية: تُبرهن المتطابقات بين المتسلسلات المصفوفية تمامًا كما
تُبرهن المتطابقات العددية، متى صدّق \omterm{def:b2:nvs:norm}{معيار} ضربي تحتي التقارب المطلق.
\end{example}

\begin{example}[معيار النهاية العليا يعني الانتظام: المعجم]\label{ex:b2:nvs:uniformdictionary}
القول $\norm{f_n - f}_\infty \to 0$ \emph{هو} التقارب
المنتظم: فعدد واحد، وهو $\sup_x\abs{f_n(x) - f(x)}$، يحدّ
الخطأ عند كل نقطة في آن واحد. والمعجم في
العمل على $f_n(x) = x^n$ فوق $\intcc{0}{1}$: نقطيًا، $f_n
\to 0$ على $\intco{0}{1}$ و$f_n(1) = 1$؛ وبالمعيار،
$\norm{f_n - 0}_\infty = 1 \not\to 0$، وفعلًا فالنهاية النقطية
غير متصلة، ومن ثم فهي بعيدة عن متناول نهاية \omterm{def:b2:nvs:norm}{بالمعيار}
$\norm\cdot_\infty$ في $C(\intcc01)$ (المغلق
تحت النهايات المنتظمة، \cref{thm:b2:metric:rncomplete}). وعلى
$\intcc{0}{a}$، $a < 1$: $\norm{f_n}_\infty = a^n \to 0$ ---
فيُستعاد التقارب المنتظم بتصغير مجال التعريف. وكل
قول في التقارب في \cref{ch:b2:funcseq} قول
عن هذا \omterm{def:b2:nvs:norm}{المعيار} الواحد؛ وإبقاء المعجم في الذهن ينصّف
ذلك الفصل.
\end{example}

\begin{remark}[مزالق شائعة]\label{rem:b2:nvs:pitfalls}
(1) يتعلق \omterm{thm:b2:nvs:continuouslinear}{معيار المؤثر} \emph{بكلا} المعيارين المختارين: فالمصفوفة
نفسها لها $\vertiii\cdot_\infty$ معطى بمجاميع الأسطر
(\cref{exo:b2:nvs:4}) و$\vertiii\cdot_1$ مختلف (مجاميع
الأعمدة)؛ وذكرُ ``\omterm{def:b2:nvs:norm}{معيار}'' مصفوفة دون تسمية
المعيارين الأساسيين لا معنى له. (2) المتراجحة $\vertiii{AB} \leq
\vertiii A\,\vertiii B$ عادةً تامة ---
فقد تتقلص القوى أسرع بكثير مما يوحي به الحدّ $\vertiii A^k$،
وهذا هو كل المغزى من المعايير المكيَّفة (مسألة نهاية
الأسبوع في هذا الفصل، السؤال 22). (3) قول ``الخطية تستلزم
الاتصال'' امتياز للبعد المنتهي: فالاشتقاق
على كثيرات الحدود خطي وغير محدود
(\cref{ex:b2:nvs:operatornorms}). (4) لا يفيد التقارب المطلق
للمتسلسلة $\sum u_n$ إلا حين يكون الفضاء \omterm{def:b2:metric:complete}{تامًا}
(ويبني \cref{exo:b2:series:9} المثالَ المضاد). (5) في
البعد غير المنتهي تكون النهاية العليا على كرة الوحدة نهاية عليا
حقيقية: فلا تفترض أنها مبلوغة
(\cref{exo:b2:nvs:8}).
\end{remark}

\begin{remark}[آفاق داخل هذا المجلد]\label{rem:b2:nvs:perspectives}
ثلاثة مواعيد صارت الآن محدَّدة. مع \cref{ch:b2:series}: في
\omterm{def:b2:nvs:banach}{فضاء باناخ}، تتقارب المتسلسلات المتقاربة بإطلاق، ومن ثم تصير
المتسلسلتان الهندسية والأسّية للمؤثرات أدواتٍ يومية ---
لقلب $I - A$، ولتعريف $\eu^{A}$
(\cref{ex:b2:series:neumann}). ومع
\cref{ch:b2:funcseq} و\cref{ch:b2:powerseries}: يكون تقارب
متتاليات الدوال ومتسلسلات القوى تقاربًا في
$\bigl(C, \norm\cdot_\infty\bigr)$
(\cref{ex:b2:nvs:uniformdictionary})، ويكون نصف قطر
التقارب قولًا في أي المتسلسلات الهندسية
تهيمن. ومع \cref{ch:b2:fourier}: يختلف المعياران
$\norm\cdot_2$ و$\norm\cdot_\infty$ اختلافًا حقيقيًا على
$C(\intcc{0}{1})$ (\cref{ex:b2:nvs:sqrtnxn})، وهذا بالضبط
هو السبب في أن تقارب متسلسلات فورييه في المتوسط التربيعي والتقارب
المنتظم مبرهنتان مختلفتان بثمنين مختلفين.
\end{remark}

\begin{remark}[أين يُستعمَل هذا الفصل]
تُدير معايير المؤثرات والمتسلسلة الهندسية حجج الاضطراب
في \cref{ch:b2:diffcalc} (مبرهنة الدالة العكسية) و
\omterm{ex:b2:nvs:matrixexp}{الأسّي المصفوفي} في \cref{ch:b2:diffeq}؛ ويرخّص \omterm{thm:b2:nvs:finitedim}{تكافؤ
المعايير} صامتًا كل حجة من نمط ``اختر معيارك المفضّل''
في \cref{ch:b2:funcseq} وما بعده؛ والفاصل بين
البعد المنتهي وغير المنتهي في مبرهنة ريس --- المُصاغ كمّيًا في
مسألة نهاية الأسبوع في هذا الفصل --- هو السبب في أن مجلد السنة
الثالثة يحتاج أدوات جديدة (التقارب الضعيف، وأرزيلا--أسكولي،
وإسقاطات فضاء هيلبرت) حيث كان هذا المجلد ما يزال يستخرج
متتاليات جزئية متقاربة.
\end{remark}

\section{تمارين}

\begin{exercise}[$\star$]\label{exo:b2:nvs:1}
على $\R^2$، ارسم كرات الوحدة للمعايير $\norm\cdot_1$ و$\norm\cdot_2$
و$\norm\cdot_\infty$، وبرهن على المتراجحات $\norm x_\infty
\leq \norm x_2 \leq \norm x_1 \leq 2\norm x_\infty$ بأفضل
الثوابت في البعد $2$.
\end{exercise}

\begin{exercise}[$\star$]\label{exo:b2:nvs:2}
هل $N(f) = \abs{f(0)} + \norm{f'}_\infty$ \omterm{def:b2:nvs:norm}{معيار} على
$C^1(\intcc{0}{1})$؟ قارنه \omterm{def:b2:nvs:norm}{بالمعيار} $\norm{f}_\infty$: فمتراجحة واحدة
تصحّ والأخرى تفشل (أبرِز ذلك).
\end{exercise}

\begin{exercise}[$\star$]\label{exo:b2:nvs:3}
احسب \omterm{thm:b2:nvs:continuouslinear}{معيار المؤثر} $u(f) = \int_0^1 f(t)\,\eu^t\,\dd t$
على $\bigl(C(\intcc{0}{1}), \norm\cdot_\infty\bigr) \to \R$، \omterm{def:b2:nvs:norm}{ومعيار}
الإزاحة $S(x_1, x_2, \dots, x_n) = (x_2, \dots, x_n, 0)$ على
$(K^n, \norm\cdot_\infty)$.
\end{exercise}

\begin{exercise}[$\star\star$]\label{exo:b2:nvs:4}
على $(\R^n, \norm\cdot_\infty)$، برهن على أن \omterm{def:b2:nvs:norm}{معيار} مؤثر
مصفوفة $A$ هو $\vertiii A_\infty = \max_i \sum_j \abs{a_{ij}}$ (أي أكبر
مجموع مطلق لسطر). واحسبه من أجل $\begin{pmatrix} 1 & -2\\
3 & 1\end{pmatrix}$.
\end{exercise}

\begin{exercise}[$\star\star$]\label{exo:b2:nvs:5}
برهن على أن $GL_n(K)$ \omterm{def:b2:metric:topology}{مفتوحة} في $\mathcal{M}_n(K)$ وأن $A
\mapsto A^{-1}$ \omterm{def:b2:metric:continuity}{متصل} عليها. \emph{إرشاد: من أجل الانفتاح، إذا كانت
$\vertiii H < \frac{1}{\vertiii{A^{-1}}}$ فإن $A + H = A(I +
A^{-1}H)$ مع $\vertiii{A^{-1}H} < 1$، وتكون $I + B$ قابلة للقلب
من أجل $\vertiii B < 1$ بالمتسلسلة الهندسية
(\cref{thm:b2:nvs:absoluteconvergence})؛ ومن أجل الاتصال، حُدّ
$(A+H)^{-1} - A^{-1}$ باستعمال المتسلسلة نفسها.}
\end{exercise}

\begin{exercise}[$\star\star$]\label{exo:b2:nvs:6}
لتكن $\varphi$ صيغة خطية على فضاء معياري $E$. برهن على أن
$\varphi$ متصلة إذا وفقط إذا كانت $\ker\varphi$ مغلقة.
\emph{(إذا كانت $\ker\varphi$ مغلقة و$\varphi \neq 0$، فاختر $a$
مع $\varphi(a) = 1$ و$r > 0$ مع $B(a, r) \cap \ker\varphi =
\emptyset$؛ واستنتج $\abs{\varphi(h)} \leq \frac{1}{r}\norm h$ بحجة
تغيير سلّم على $a - \frac{h}{\varphi(h)}$.)}
\end{exercise}

\begin{exercise}[$\star\star$]\label{exo:b2:nvs:7}
برهن على أن $\bigl(C(\intcc{0}{1}), \norm\cdot_1\bigr)$ ليس
\omterm{def:b2:metric:complete}{تامًا}: بيّن أن الدوال $f_n$، وهي منحدرات أفينية من $0$ إلى
$1$ على $\bigl[\frac12 - \frac1n, \frac12\bigr]$ (بالقيمة $0$
قبله و$1$ بعده)، تشكّل متتالية كوشي بلا نهاية متصلة
من أجل $\norm\cdot_1$.
\end{exercise}

\begin{exercise}[$\star\star\star$]\label{exo:b2:nvs:8}
على $E = C(\intcc{0}{1})$ مع $\norm\cdot_\infty$، لننظر في
\[
\varphi(f) = \sum_{n \geq 1} (-1)^n\, 2^{-n} f\bigl(\tfrac1n\bigr).
\]
برهن على أن $\varphi$ صيغة خطية متصلة معرَّفة جيدًا مع
$\vertiii\varphi = 1$، لكن النهاية العليا التي تعرّف
$\vertiii\varphi$ \emph{ليست مبلوغة} على كرة الوحدة المغلقة.
\emph{(الحدّ الأعلى: بمتراجحة المثلث. \omterm{def:b2:nvs:norm}{والمعيار} $= 1$: ابنِ
دوالًا متصلة $f_K$ بحيث $\norm{f_K}_\infty \leq 1$ و$f_K(\frac1n)
= (-1)^n$ من أجل $n \leq K$ --- فالنقاط $\frac1n$ معزولة
بعضها عن بعض. وعدم البلوغ: تفرض المساواة أن $f(\frac1n) =
(-1)^n$ من أجل كل $n$، وهذا يتنافى مع اتصال $f$ عند $0$
لأن $\frac1n \to 0$.)}
\end{exercise}

\begin{exercise}[$\star\star\star$]\label{exo:b2:nvs:9}
ليكن $E$ فضاءً معياريًا كرة الوحدة المغلقة فيه متراصة.
أعد استنباط أن لكل متتالية محدودة متتاليةً جزئية متقاربة، دون
الاستشهاد بالمبرهنة \cref{thm:b2:nvs:riesz}، وبرهن على أن كل
صيغة خطية على $E$ متصلة إذا وفقط إذا كان $\dim E < \infty$.
\emph{(من أجل البعد غير المنتهي، ابنِ صيغة غير متصلة
بتعريفها بحرية على متتالية مستقلة خطيًا منظَّمة
ثم تمديدها --- مع قبول وجود مكمّل
جبري.)}
\end{exercise}

\begin{exercise}[$\star\star$]\label{exo:b2:nvs:10}
على $C(\intcc{0}{1})$، برهن على أن $\norm f_1 \leq \norm f_2 \leq \norm
f_\infty$ \emph{(بكوشي--شوارتز من أجل الأولى)}، وبيّن بالعائلة
$f_n(x) = x^n$ أنه لا يمكن عكس أي من المتراجحتين
إلى حدّ ثابت: أي إن المعايير الثلاثة غير متكافئة مثنى مثنى.
\end{exercise}

\begin{exercise}[$\star\star$]\label{exo:b2:nvs:11}
(المسافة إلى فضاء فوقي) لتكن $\varphi$ صيغة خطية متصلة غير
معدومة على فضاء معياري $E$. برهن على أن
\[
d\bigl(x, \ker\varphi\bigr) =
\frac{\abs{\varphi(x)}}{\vertiii\varphi}
\qquad (x \in E),
\]
وتحقق على \cref{exo:b2:nvs:8} من أن النهاية الدنيا لا يلزم أن
تُبلغ عند أي نقطة من الفضاء الفوقي.
\end{exercise}

\begin{exercise}[$\star\star\star$]\label{exo:b2:nvs:12}
على $E = \R[X]$ (أي كل كثيرات الحدود)، ليكن $N_1(P) =
\sup_{\intcc{0}{1}}\abs P$ و$N_2(P) =
\sup_{\intcc{0}{2}}\abs P$. بيّن أن $N_1 \leq N_2$ لكن
$N_1$ و$N_2$ \emph{غير} متكافئين؛ واستنتج أن
المطابق $(E, N_2) \to (E, N_1)$ تقابل خطي \omterm{def:b2:metric:continuity}{متصل}
مقلوبه غير \omterm{def:b2:metric:continuity}{متصل}. وبيّن أخيرًا أن $(E, N_1)$
ليس \omterm{def:b2:metric:complete}{تامًا} \emph{(بالمجاميع الجزئية لتايلور للدالة $\eu^x$)}. والظواهر الثلاث
مستحيلة في البعد المنتهي --- فقل لماذا.
\end{exercise}

\section{مسألة: أفضل تقريب ومبرهنة تشيبيشيف}

إلى أي حدّ يمكن تقريب دالة بكثيرات حدود من درجة معطاة،
وأي كثير حدود يقرّبها أفضل تقريب؟ أما جانب الوجود
فجوابه من هذا الفصل: فالتراص في البعد المنتهي
يجعل أفضل التقريبات موجودة. وأما الجانب الصريح،
فحالة واحدة غير بديهية تُحلّ كاملةً باليدين المجرّدتين ---
فمن بين كل كثيرات الحدود \emph{الموحَّدة} من الدرجة $n$، أصغرها
معيارًا أعلى على $\intcc{-1}{1}$ هو كثير حدود تشيبيشيف (المنظَّم)،
ومعياره $2^{1-n}$: وهي \emph{مبرهنة تشيبيشيف
الحدّية}. وتبرهن المسألة على الجانبين، ثم تقيس إلى أي حدّ
يفشل التراص في البعد غير المنتهي: فكرة الوحدة في
$C(\intcc{0}{1})$ تحتوي كوكبات غير منتهية من النقاط على
مسافة متبادلة $1$.

\begin{problem}[{مسألة نهاية الأسبوع --- مبرهنة تشيبيشيف الحدّية
وهندسة كرة الوحدة}]
\label{pb:b2:nvs:1}
المعايير بلا دليل سفلي هي معايير النهاية العليا على القطعة المذكورة.

\medskip
\noindent\textbf{الجزء الأول --- أفضل تقريب في \omterm{def:b2:nvs:norm}{الفضاءات
المعيارية}.}
\begin{enumerate}
  \item ليكن $F$ فضاءً جزئيًا منتهي البعد من فضاء معياري
        $E$ وليكن $x \in E$. برهن على أن المسافة $d(x,
        F) = \inf_{f \in F}\norm{x - f}$ \emph{مبلوغة}
        \emph{(أرجِع المسألة إلى جزء مغلق محدود من $F$ واستعمل
        \cref{thm:b2:nvs:finitedim})}.
  \item يكون \omterm{def:b2:nvs:norm}{المعيار} \emph{محدبًا تمامًا} إذا كان $\norm u = \norm v
        = 1$ و$u \neq v$ يستلزمان $\bigl\Vert\frac{u +
        v}2\bigr\Vert < 1$. بيّن أن $\norm\cdot_2$ على $\R^n$
        محدب تمامًا \emph{(بمتطابقة متوازي الأضلاع)}، وأن
        $\norm\cdot_1$ و$\norm\cdot_\infty$ ليسا كذلك من أجل
        $n \geq 2$.
  \item برهن على أن أفضل تقريب في السؤال 1 \emph{وحيد} من أجل
        \omterm{def:b2:nvs:norm}{معيار} محدب تمامًا.
  \item في $(\R^2, \norm\cdot_\infty)$، احسب كل أفضل
        التقريبات للعنصر $x = (0, 1)$ بالمستقيم $F =
        \operatorname{Vect}\bigl((1,0)\bigr)$: فتحصل على فترة
        من المصغِّرات.
  \item في $\bigl(C(\intcc{a}{b}), \norm\cdot_\infty\bigr)$،
        بيّن أن أفضل تقريب للدالة $f$ \emph{بالثوابت}
        وحيد ويساوي $c^* = \frac{\max f
        + \min f}{2}$، بمسافة $\frac{\max f - \min
        f}{2}$؛ واحسب الاثنين من أجل $f(x) = x^2$ على
        $\intcc{0}{1}$.
\end{enumerate}

\medskip
\noindent\textbf{الجزء الثاني --- كثيرات حدود تشيبيشيف.}
\begin{enumerate}[resume]
  \item بيّن أنه يوجد كثير حدود واحد بالضبط $T_n$ يحقق
        $T_n(\cos\theta) = \cos n\theta$ من أجل كل $\theta$
        \emph{(بالتراجع $T_{n+1} = 2XT_n - T_{n-1}$ الآتي من
        صيغة جمع جيوب التمام)}، وأن $\deg T_n = n$، وأن
        معامله الرئيسي $2^{n-1}$ من أجل $n \geq 1$.
  \item بيّن أن $\abs{T_n} \leq 1$ على $\intcc{-1}{1}$، مع
        $T_n(\eta_k) = (-1)^k$ عند النقاط $n + 1$ وعددها $\eta_k =
        \cos\frac{k\pi}{n}$ (مع $k = 0, \dots, n$)، وأن
        جذور $T_n$ هي النقاط $n$ وعددها
        $\cos\frac{(2k-1)\pi}{2n}$، متداخلةً مع النقاط $\eta_k$.
  \item احسب $T_2, T_3, T_4$، وتحقق من تناوب
        $T_3$ عند $\eta_0, \dots, \eta_3 = 1, \frac12, -\frac12,
        -1$ بالتقييم المباشر.
  \item من أجل $\abs x \geq 1$، برهن على أن
        \[
        T_n(x) = \frac{\bigl(x + \sqrt{x^2 - 1}\bigr)^n +
        \bigl(x - \sqrt{x^2 - 1}\bigr)^n}{2},
        \]
        واستنتج $T_n(x) \sim \frac12\bigl(x + \sqrt{x^2 -
        1}\bigr)^n \to \infty$ هندسيًا من أجل $x > 1$ ثابت.
  \item برهن على قانون التركيب $T_m \circ T_n = T_{mn}$
        \emph{(تحقق منه على $\intcc{-1}{1}$ واستشهد بصلابة
        كثيرات الحدود)}.
\end{enumerate}

\medskip
\noindent\textbf{الجزء الثالث --- مبرهنة تشيبيشيف الحدّية.}
نكتب $Q_n = 2^{1-n}T_n$ (وهو موحَّد، حسب السؤال 6).
\begin{enumerate}[resume]
  \item ليكن $P$ موحَّدًا من الدرجة $n \geq 1$ مع
        $\sup_{\intcc{-1}{1}}\abs P < 2^{1-n}$. وبتقييم $D
        = Q_n - P$ عند النقاط $\eta_k$ وعدّ تغيّرات
        الإشارة، استنتج تناقضًا. واختم بأن:
        \[
        \sup_{\intcc{-1}{1}}\abs P \;\geq\; 2^{1-n}
        \qquad\text{من أجل كل كثير حدود موحَّد } P \text{ من الدرجة } n.
        \]
  \item (حالة المساواة) لنفترض $\sup_{\intcc{-1}{1}}\abs P =
        2^{1-n}$ مع $P$ موحَّد من الدرجة $n$، وليكن $D = Q_n
        - P \neq 0$. بيّن أن $(-1)^kD(\eta_k) \geq 0$ من أجل كل $k$؛
        وبيّن أن كلًا من الفترات $n$
        $\intcc{\eta_{k}}{\eta_{k-1}}$ يحتوي جذرًا للكثير $D$،
        وأن الجذر المشترك بين فترتين متتاليتين نقطةٌ
        \emph{داخلية} $\eta_k$ يكون عندها $D' = 0$ أيضًا. واختم بأن
        للكثير $D$ عدد $n$ من الجذور محسوبةً بتضاعفاتها، ومن ثم $D = 0$: أي إن المصغِّر هو بالضبط
        $Q_n$ --- وهي \emph{مبرهنة تشيبيشيف الحدّية}.
  \item أعد صياغة المبرهنة بوصفها مسافة: على $\intcc{-1}{1}$،
        \[
        d_\infty\bigl(X^n,\ \R_{n-1}[X]\bigr) = 2^{1-n},
        \]
        بأفضل تقريب وحيد $X^n - Q_n$؛ وبيّن بالتعويض الأفيني
        $x = \frac{1+t}2$ أن المسافة تصير على
        $\intcc{0}{1}$ هي $2^{1-2n}$.
  \item (عقد الاستكمال المثلى) من أجل $n$ عقدة $x_1, \dots,
        x_n \in \intcc{-1}{1}$، يكون كثير حدود العقد
        $\omega(x) = \prod_i(x - x_i)$ موحَّدًا من الدرجة $n$.
        استنتج من السؤال 12 أي اختيار للعقد يصغّر
        $\sup_{\intcc{-1}{1}}\abs\omega$، وهو العامل المتعلق بالعقد
        في حدّ خطأ الاستكمال الكلاسيكي،
        وأعطِ القيمة الدنيا.
  \item تحقق يدويًا من حالة $n = 2$ في المبرهنة (جد
        $\inf_c \sup_{\intcc{-1}{1}}\abs{x^2 - c}$ مباشرةً)،
        واحسب عدديًا مسافة السؤال 13 على
        $\intcc{0}{1}$ من أجل $n = 10$. وماذا يقول حجمها
        عن منحني $x^{10}$؟
\end{enumerate}

\medskip
\noindent\textbf{الجزء الرابع --- كرة الوحدة في
$C(\intcc{0}{1})$.}
\begin{enumerate}[resume]
  \item لتكن $g_k(x) = x^{2^k}$. بيّن أن $\norm{g_k}_\infty = 1$ وأن
        $\norm{g_k - g_j}_\infty \geq \frac14$ من أجل $j > k$
        \emph{(قيّم عند النقطة التي يكون فيها $x^{2^k} =
        \frac12$)}: فنحصل على متتالية محدودة صريحة بلا
        متتالية جزئية متقاربة --- أي إن كرة الوحدة المغلقة غير
        متراصة، باليدين المجرّدتين.
  \item (مبرهنة ريس المساعدة، مشحوذةً) ليكن $F$ فضاءً جزئيًا فعليًا
        \emph{منتهي البعد} من فضاء معياري
        $E$. باستعمال السؤال 1، أنتج متجهة واحدية $x$
        تحقق $d(x, F) = 1$ بالضبط --- لا مجرد $\geq 1 -
        \varepsilon$ كما في مبرهنة \cref{thm:b2:nvs:riesz} المساعدة.
  \item استنتج: في كل فضاء معياري غير منتهي البعد توجد
        متتالية من متجهات الوحدة المسافات بينها مثنى مثنى
        $\geq 1$، وأعد استنباط مبرهنة ريس منها.
  \item في $C(\intcc{0}{1})$، أبرِز كوكبة كهذه
        صراحةً: أي دوال الخيمة $h_n$ ذوات الحامل
        $\bigl[\frac1{n+1}, \frac1n\bigr]$ والقيمة القصوى $1$.
        وتحقق من $\norm{h_n} = 1$، ومن $\norm{h_n - h_m} = 1$ من أجل $n
        \neq m$، ولاحظ أن $h_n \to 0$ نقطيًا لا
        بانتظام.
  \item (يفشل الحدّ الكلي) بيّن أن كرة الوحدة المغلقة
        في $C(\intcc{0}{1})$ لا يمكن تغطيتها بعدد منته من
        الكرات ذوات نصف القطر $\frac13$ \emph{(فكل كرة كهذه تحتوي
        عنصرًا واحدًا على الأكثر من العناصر $h_n$)} --- وقارِن بخطوة
        الحدّ الكلي في برهان
        \cref{thm:b2:metric:borellebesgue}.
\end{enumerate}

\medskip
\noindent\textbf{الجزء الخامس --- المعايير في العمل على المصفوفات،
والتركيب.}
\begin{enumerate}[resume]
  \item برهن على أن كل \omterm{def:b2:reduction:eigen}{قيمة ذاتية} $\lambda$ للمصفوفة $A \in
        \mathcal{M}_n(\C)$ تحقق $\abs\lambda \leq
        \vertiii A$ من أجل كل \omterm{def:b2:nvs:norm}{معيار} مؤثر؛ وطبّق
        \cref{exo:b2:nvs:4} لتحدّ \omterm{def:b2:reduction:eigen}{القيم الذاتية} للمصفوفة
        $\begin{psmallmatrix}1 & -2\\ 3 & 1\end{psmallmatrix}$
        وقارِن بطويلتها الحقيقية.
  \item (المعايير المكيَّفة) لتكن $A$ قابلة للتقطير، مع $A =
        P\,\mathrm{diag}(\lambda_1, \dots, \lambda_n)\,P^{-1}$.
        بيّن أن $N_P(x) = \norm{P^{-1}x}_\infty$ \omterm{def:b2:nvs:norm}{معيار} يحقق \omterm{def:b2:nvs:norm}{معيار}
        مؤثره $\vertiii A_{N_P} =
        \max_i\abs{\lambda_i}$.
  \item استنتج: من أجل $A$ قابلة للتقطير، يكون $A^k \to 0$ إذا وفقط
        إذا حققت كل \omterm{def:b2:reduction:eigen}{القيم الذاتية} $\abs{\lambda_i} < 1$ ---
        ويجعل \omterm{thm:b2:nvs:finitedim}{تكافؤ المعايير} النتيجةَ
        غير متعلقة \omterm{def:b2:nvs:norm}{بالمعيار}. وتحقق على $A =
        \frac14\begin{psmallmatrix}1 & 2\\ 2 &
        1\end{psmallmatrix}$.
  \item (ثوابت التكافؤ تنفجر) على $\R_n[X]$، قارِن
        $N_c(P) = \max_k \abs{a_k}$ (بالمعاملات) و
        $\norm{P}_{\intcc{0}{1}}$: فكلاهما \omterm{def:b2:nvs:norm}{معيار}، ومن ثم فهما
        متكافئان من أجل كل $n$ ثابت؛ لكن بيّن، باستعمال المصغِّر
        الموحَّد في السؤال 13 على $\intcc{0}{1}$، أن
        أفضل ثابت $C_n$ في $N_c \leq
        C_n\norm\cdot_{\intcc{0}{1}}$ يحقق $C_n \geq
        2^{2n-1}$. واختم بجملة واحدة عن سبب موت قول ``كل \omterm{def:b2:nvs:equivalent}{المعايير
        متكافئة}'' في البعد غير المنتهي.
  \item (تركيب) بجملة واحدة لكل بند: أين عمل تراص
        كرات البعد المنتهي (السؤالان 1 و12)؛ وماذا يحكم
        التحدب التام؛ وماذا تهدم كوكبة
        السؤالين 18--19؛ وكيف يقيس السؤال 24
        الفشل كمّيًا. وسمِّ القمة (مبرهنة تشيبيشيف
        الحدّية) وقل أين يجد أفضل تقريب موطنه
        الحديث (مبرهنة الإسقاط على فضاءات هيلبرت،
        مجلد السنة الثالثة، حيث يحلّ التمام محلّ
        التراص).
\end{enumerate}
\end{problem}
